\chapter{Résultats et validation}
\label{chap:resultats}

\section{Synthèse}

\begin{table}[H]
\centering
\caption{Indicateurs de la résolution sur le scénario de référence}
\label{tab:synthese-resultats}
\begin{tabular}{@{}ll@{}}
\toprule
\textbf{Indicateur} & \textbf{Valeur} \\
\midrule
Statut du solveur & \textbf{Optimal} \\
Temps de résolution & \SI{0,06}{\second} (trois passes) \\
Taille du problème & 170 variables dont 29 binaires, 138 contraintes \\
Demande satisfaite & \textbf{\SI{100}{\percent}} --- aucun manque \\
Validation indépendante & \textbf{86 contrôles réussis sur 86} \\
Niveau 1 (demande non servie) & \num{0,000} \\
Niveau 2 (violations molles) & \num{1268,3} \\
Niveau 3 (coût opératoire) & \num{3111,9} \\
\bottomrule
\end{tabular}
\end{table}

Toute la demande est servie et la solution respecte l'intégralité des règles de procédé.
Mais le système \emph{ne tient pas dans ses limites de stockage} : deux violations
structurelles subsistent, qu'aucune décision d'exploitation ne peut corriger.

\section{Le plan d'exploitation optimal}

\subsection{Décadmiation}

Une seule ligne décadmie, au niveau minimal non nul : \SI{750}{\tonne} sur 13XY, soit un
filtre. C'est le strict nécessaire pour couvrir le besoin de U116BC en acide 29 décadmié
(\SI{363,6}{\tonne}) et alimenter la production obligatoire de DEC\_CL.

Le niveau 3 de l'objectif pénalisant la décadmiation, l'optimiseur ne fait rien de plus que
ce que les contraintes exigent.

\subsection{Transferts interzones}

\begin{table}[H]
\centering
\caption{Transferts retenus par l'optimisation}
\label{tab:transferts}
\begin{tabular}{@{}llS[table-format=3.1]@{}}
\toprule
\textbf{Origine} & \textbf{Destination} & {\textbf{Quantité} (\tP)} \\
\midrule
13F & 13ZU & 758.8 \\
13F & 13E & 603.3 \\
13CD & 13AB & 100.0 \\
\bottomrule
\end{tabular}
\end{table}

Trois transferts seulement, et leur logique est limpide. La ligne 13F est la source : seule
ligne sans concentration associée, elle produit \SI{1500}{\tonne} qu'elle n'a presque pas
moyen de consommer. Les lignes 13ZU et 13E sont les puits : la première ne produit que
\SI{800}{\tonne} mais doit charger \SI{1481,7}{\tonne} en concentration, la seconde produit
\SI{1500}{\tonne} mais doit charger \SI{1505}{\tonne} \emph{et} livrer \SI{700}{\tonne} à
IMACID. Le transfert 13CD $\to$ 13AB est exactement au minimum d'expédition : c'est un
appoint, non un flux de fond.

\begin{resultat}[title={Un indice de validité}]
Le modèle retrouve seul le rôle que la ligne 13F joue réellement dans l'usine : celui de
\textbf{réserve souple} du système. Cette concordance entre un résultat calculé et une
pratique d'exploitation constitue un bon signe de validité de la formulation.
\end{resultat}

\subsection{Clarification}

\begin{table}[H]
\centering
\caption{Répartition de la charge des décanteurs}
\label{tab:clarification}
\begin{tabular}{@{}lS[table-format=3.1]S[table-format=3.1]S[table-format=3.1]S[table-format=2.1]@{}}
\toprule
\textbf{Ligne} & {\textbf{NCL}} & {\textbf{DEC}} & {\textbf{Total}} & {\textbf{Occupation} (\si{\percent})} \\
\midrule
14XY & 0.0 & 476.7 & 476.7 & 47.7 \\
14ZU & 722.5 & 0.0 & 722.5 & 72.3 \\
\bottomrule
\end{tabular}
\end{table}

La séparation est nette : 14XY, seule ligne habilitée à produire du DEC\_CL, y consacre ses
décanteurs ; 14ZU produit l'acide clarifié destiné à IR12, donc à U53. Les lignes 14AB et
14CD, bien qu'habilitées, ne clarifient pas --- ce n'est pas nécessaire, IR12 contenant déjà
\SI{3821,6}{\tonne} pour une demande de \SI{1500}{\tonne}.

\section{Les deux violations structurelles}

\subsection{IR11 déborde de 1\,016 tonnes par jour}

\begin{equation}
\underbrace{\num{5514,7}}_{\text{initial}} + \underbrace{\num{2860,4}}_{\text{entrées}} - \underbrace{\num{986,7}}_{\text{sorties}} = \num{7388,3}\tP
\qquad\text{pour une capacité de } \num{6372,4}\tP
\end{equation}

soit un remplissage de \SI{115,9}{\percent}.

\paragraph{Origine.} La cocristallisation est systématique : les quatre échelons de 14EXT et
quatre des six échelons de 14AB y sont affectés et traitent donc tout ce qu'ils produisent,
quelle que soit la demande. Il en résulte \SI{2431}{\tonne} de CoC par jour pour un besoin
en acide de qualité décadmiée de \SI{987}{\tonne} seulement.

\paragraph{Pourquoi l'optimiseur ne peut rien y faire.} La production de CoC est un
\emph{paramètre}, non une variable ; les sorties sont bornées par la demande, qui est fixée.
Le solde net atteint \SI{+1444}{\tonne} par jour, dans un bac qui n'offrait que
\SI{858}{\tonne} de marge.

\subsection{Le stock de 14EXT reste sous sa bande de sécurité}

Le stock final s'établit à \SI{220,8}{\tonne} pour une bande de
\SIrange{473,1}{1482,4}{\tonne}, soit un manque de \SI{252,3}{\tonne}.

\paragraph{Origine.} Les \emph{quatre} échelons de 14EXT étant affectés à la
cocristallisation, la ligne ne produit \textbf{aucun} acide 54 NCL ordinaire. Son stock ne
peut donc que rester à son niveau initial, déjà trop bas.

\subsection{Une cause unique}

Les deux violations remontent à la même décision de configuration : le nombre d'échelons
affectés à la cocristallisation.

\section{Analyse de sensibilité}

\subsection{Nombre d'échelons de 14EXT en cocristallisation}

\begin{table}[H]
\centering
\caption{Effet du nombre d'échelons cocristallisants de 14EXT}
\label{tab:sensib-coc}
\begin{tabular}{@{}lS[table-format=4.1]S[table-format=4.1]S[table-format=4.1]S[table-format=4.1]@{}}
\toprule
\textbf{Configuration} & {\textbf{CoC produit}} & {\textbf{Stock IR11}} & {\textbf{Dépassement}} & {\textbf{Niveau 2}} \\
\midrule
4 sur 4 \emph{(actuel)} & 2431.3 & 7388.3 & 1016.0 & 1268.3 \\
3 sur 4 & 2235.3 & 7157.7 & 785.4 & 792.7 \\
2 sur 4 & 1899.3 & 6762.5 & 390.1 & 390.1 \\
\textbf{1 sur 4} & 1563.3 & 6367.2 & \bfseries 0.0 & \bfseries 0.0 \\
0 sur 4 & 1227.3 & 5971.9 & 0.0 & 0.0 \\
\bottomrule
\end{tabular}
\end{table}

\begin{resultat}[title={Recommandation principale}]
Ramener de \textbf{quatre à un} le nombre d'échelons de 14EXT affectés à la cocristallisation
résorbe \textbf{intégralement} les deux violations --- le niveau 2 passe de \num{1268,3} à
\num{0} --- tout en continuant à servir \SI{100}{\percent} de la demande.

Le gain est double : IR11 rentre dans sa capacité, \emph{et} la ligne 14EXT se remet à
produire de l'acide 54 NCL ordinaire, ce qui ramène son stock dans sa bande de sécurité.

Cette action ne demande \textbf{aucun investissement} : seulement une reconfiguration de la
marche des échelons.
\end{resultat}

\subsection{Part minimale de DEC\_CL dans IR11 (hypothèse H8)}

\begin{table}[H]
\centering
\caption{Effet du paramètre $\alpha$}
\label{tab:sensib-alpha}
\begin{tabular}{@{}S[table-format=1.2]S[table-format=3.1]S[table-format=4.1]S[table-format=4.1]l@{}}
\toprule
{$\bm{\alpha}$} & {\textbf{DEC\_CL}} & {\textbf{Stock IR11}} & {\textbf{Dépassement}} & \textbf{Statut} \\
\midrule
0.00 & 0.0 & 6959.3 & 586.9 & optimal \\
0.05 & 128.0 & 7087.2 & 714.9 & optimal \\
0.10 & 270.1 & 7229.4 & 857.1 & optimal \\
0.15 & 429.1 & 7388.3 & 1016.0 & optimal \emph{(hypothèse)} \\
0.20 & 607.8 & 7567.1 & 1194.8 & optimal \\
0.25 & {---} & {---} & {---} & \textbf{infaisable} \\
0.30 & {---} & {---} & {---} & \textbf{infaisable} \\
\bottomrule
\end{tabular}
\end{table}

Deux enseignements.

\paragraph{Le domaine admissible de $\alpha$ est borné : $\alpha \lesssim 0{,}22$.}
Au-delà, le modèle n'a plus de solution. La raison est structurelle : \emph{seule} la ligne
14XY est habilitée à produire du DEC\_CL, et la capacité de ses échelons non
cocristallisants (\SI{820,8}{\tonne}) plafonne la production possible.

\begin{encadre}[title={Information à remonter à l'encadrant}]
Si la valeur métier de $\alpha$ dépasse \num{0,22}, la configuration actuelle ne peut pas la
respecter. Il faudrait habiliter une \textbf{seconde ligne} à produire du DEC\_CL. Cette
conclusion a pu être établie \emph{sans connaître} la valeur de $\alpha$.
\end{encadre}

\paragraph{La contrainte de qualité n'est pas la cause du débordement.} Même avec
$\alpha = 0$, IR11 déborde encore de \SI{586,9}{\tonne}. La contrainte ne fait qu'accentuer
un problème dont l'origine est ailleurs.

\subsection{Borne maximale des transferts (hypothèse H7)}

\begin{table}[H]
\centering
\caption{Effet du paramètre $\tau^{\max}$}
\label{tab:sensib-tau}
\begin{tabular}{@{}S[table-format=4.0]S[table-format=1.0]S[table-format=4.1]S[table-format=4.1]@{}}
\toprule
{$\bm{\tau^{\max}}$ (\tP)} & {\textbf{Nombre de transferts}} & {\textbf{Niveau 2}} & {\textbf{Niveau 3}} \\
\midrule
300 & 6 & 1430.4 & 3392.3 \\
500 & 4 & 1268.3 & 3161.9 \\
750 & 3 & 1268.3 & 3111.9 \\
1000 & 3 & 1268.3 & 3111.9 \\
1500 & 3 & 1268.3 & 3111.9 \\
\bottomrule
\end{tabular}
\end{table}

\begin{resultat}[title={Une incertitude transformée en non-problème}]
À partir de \SI{750}{\tonne}, la solution ne bouge plus : le choix de \SI{1000}{\tonne}
comme valeur d'attente n'a \emph{aucune} influence sur le résultat.

L'incertitude de l'hypothèse H7 est donc sans conséquence. La question posée à l'encadrant
passe de « priorité haute » à simplement informative --- \textbf{elle a été résolue sans
qu'une réponse soit nécessaire}.

C'est précisément ce qu'on attend d'une analyse de sensibilité : elle ne devine pas la
donnée manquante, elle mesure l'effet de son absence.
\end{resultat}

\section{Goulots d'étranglement}

\begin{table}[H]
\centering
\caption{Ressources les plus tendues à l'optimum}
\label{tab:goulots}
\begin{tabular}{@{}lS[table-format=4.1]S[table-format=4.1]S[table-format=3.1]l@{}}
\toprule
\textbf{Ressource} & {\textbf{Utilisé}} & {\textbf{Capacité}} & {\textbf{Taux} (\si{\percent})} & \\
\midrule
Bac IR11 & 7388.3 & 6372.4 & 115.9 & dépassé \\
Cuves 13F & 965.3 & 965.3 & 100.0 & saturé \\
Cuves 13XY & 965.3 & 965.3 & 100.0 & saturé \\
Cuves 13ZU & 965.3 & 965.3 & 100.0 & saturé \\
Cuves 13CD & 950.7 & 965.3 & 98.5 & \\
Décanteurs 14ZU & 722.5 & 1000.0 & 72.3 & \\
Cuves 13AB & 650.2 & 965.3 & 67.4 & \\
Cuves 13E & 641.7 & 965.3 & 66.5 & \\
\bottomrule
\end{tabular}
\end{table}

\paragraph{Le stockage d'acide 29 est saturé sur trois lignes.} C'est ce qui explique
pourquoi 13F ne transfère « que » \SI{1362}{\tonne} alors qu'elle est excédentaire : il n'y
a plus de place ailleurs. Le système est contraint par le \emph{stockage}, non par la
production.

\paragraph{Les décanteurs ne sont pas un goulot dans ce scénario} (\SI{72}{\percent} au
plus). Ce constat est contre-intuitif au regard du chapitre~\ref{chap:procede}, qui les
présente comme la contrainte structurante --- et c'est un enseignement en soi : la contrainte
\emph{conceptuellement} la plus subtile n'est pas nécessairement la plus \emph{limitante} un
jour donné.

\begin{alerte}[title={Ne pas généraliser}]
Ce diagnostic vaut pour \emph{ce} scénario. Avec une demande d'acide clarifié plus forte, ou
un $\alpha$ plus élevé, les décanteurs redeviendraient contraignants. C'est précisément
l'intérêt de disposer d'un modèle : rejouer le calcul chaque jour plutôt que raisonner sur
des intuitions figées.
\end{alerte}

\section{Validation}

\subsection{Le validateur indépendant}

\num{86} contrôles, tous réussis : bilans matière fermés sur les dix-neuf n\oe uds du
système, conservation globale du \PdeuxOcinq{}, demandes exactement servies, capacités
respectées, niveaux de décadmiation discrets, rendements exacts, stocks positifs.

\subsection{La tolérance de concentration n'est pas utilisée}

L'encadrant a accordé une tolérance sur la charge de concentration tout en la qualifiant de
« normalement rigide ». L'écart constaté est nul sur les cinq lignes : la bande existe en
dernier recours, la pénalité du niveau 2 maintient le modèle à l'égalité stricte. Le
comportement obtenu est exactement conforme à la consigne.

\subsection{La contrainte de qualité d'IR11 est nécessaire}

La contre-épreuve automatisée présentée au chapitre~\ref{chap:implementation} confirme
qu'avec $\alpha = 0$, la production de DEC\_CL tombe à zéro. La décision de modélisation est
donc empiriquement justifiée, et non seulement argumentée.

\section{Synthèse pour l'exploitation}

\subsection{Ce qui fonctionne}

Le modèle sert \SI{100}{\percent} de la demande en \SI{0,06}{\second}, sa solution est
physiquement valide, et le plan produit est sobre : une seule décadmiation, trois transferts,
deux lignes clarifiantes. Rien d'inutile.

\subsection{Ce qui ne va pas, et qui ne relève pas de l'optimisation}

\begin{enumerate}
  \item IR11 déborde de \SI{1016}{\tonne} par jour --- trop d'échelons en cocristallisation ;
  \item le stock de 14EXT reste sous sa bande --- même cause ;
  \item le stockage d'acide 29 est saturé sur trois lignes sur six.
\end{enumerate}

\subsection{Trois questions posées à l'encadrant}

\begin{enumerate}
  \item Peut-on réduire à un ou deux le nombre d'échelons de 14EXT en cocristallisation ?
        C'est la seule action qui résorbe les deux violations, et elle ne coûte rien.
  \item L'excédent de CoC a-t-il un débouché non modélisé --- vente d'acide marchand, export,
        autre bac ? Si oui, le débordement d'IR11 est un artefact de modélisation qu'il
        suffit de corriger ; sinon, c'est un problème d'exploitation réel que le modèle vient
        de mettre au jour.
  \item Quelle est la vraie valeur de $\alpha$ ? Si elle dépasse \num{0,22}, il faut habiliter
        une deuxième ligne au DEC\_CL.
\end{enumerate}

\begin{encadre}[title={Le livrable véritable}]
Le résultat le plus utile de cette étude n'est pas le plan quotidien, mais le
\textbf{diagnostic} qui l'accompagne. Un modèle d'optimisation ne doit pas seulement dire
quoi faire : il doit dire ce qui empêche de faire mieux.
\end{encadre}
